\documentclass[11pt,a4paper]{article}
\usepackage[T1]{fontenc}
\usepackage{lmodern,microtype}
\usepackage[margin=25mm]{geometry}
\usepackage{amsmath,amssymb,amsthm}
\usepackage[hidelinks]{hyperref}
\usepackage{enumitem}
\setlength{\parindent}{0pt}
\setlength{\parskip}{6pt}
\setlist{nosep}
\newtheorem{theorem}{Theorem}
\newtheorem*{faltings}{Faltings' theorem (used without proof)}
\newcommand{\C}{\mathbb C}
\newcommand{\Q}{\mathbb Q}
\newcommand{\Z}{\mathbb Z}
\hypersetup{pdftitle={Finiteness of the integer solutions of x4 + xy + y3 + 1 = 0},
  pdfsubject={A correction of the requested infinitude claim, with explicit theorem dependencies}}
\begin{document}
\begin{center}
  {\Large\bfseries Finiteness of the integer solutions of}\par\vspace{5pt}
  {\Large $x^4+xy+y^3+1=0$}
\end{center}

\textbf{Answer.} The proposed infinitude statement is false. In fact, this
equation has only finitely many rational solutions, and hence only finitely
many integer solutions.

\textbf{Scope of the proof.} The algebraic
calculations below are supplied in full. The finiteness conclusion uses
Faltings' theorem, and the genus calculation uses the standard genus formula
for a smooth plane curve. These results are stated explicitly and are not
proved here. Thus the argument is unconditional, but it is not a fully
self-contained elementary proof. A proof of the requested infinitude cannot
be supplied because that conclusion is false.

\begin{theorem}
The set
\[
  \{(x,y)\in\Q^2:x^4+xy+y^3+1=0\}
\]
is finite. It contains $(0,-1)$ and $(1,-1)$.
\end{theorem}

\subsection*{1. Intuition: check the curve before seeking a family}
One equation in two variables need not have infinitely many integer
solutions. Here the natural object to examine is the plane curve itself.
The total degree is four, but that alone does not determine its arithmetic:
factorization or singularities could change the situation substantially.
We therefore check two precise facts: the curve is geometrically
irreducible, and its projective completion is smooth. A smooth plane quartic
has genus three, which places it within the scope of a theorem asserting
finiteness even for rational points.

The most useful algebraic observation is that, at a hypothetical singular
point, the derivative equations become linear relations among the three
monomials $x^4$, $xy$, and $y^3$. These relations force incompatible absolute
values. This makes the essential smoothness check particularly short.

\subsection*{2. Geometric irreducibility}
Put $f(x,y)=x^4+xy+y^3+1$. As a polynomial in $y$ over the field $\C(x)$,
it has degree three. A reducible cubic over a field has a root in that field.
Suppose such a root were $r=p/q$, where $p,q\in\C[x]$ are coprime and
$q\ne0$. Clearing denominators gives
\[
  p^3+xpq^2+(x^4+1)q^3=0.
\]
Thus $q$ divides $p^3$. Coprimality forces $q$ to be constant, so $r$ is a
polynomial. If $\deg r\le1$, the term $x^4$ has strictly greater degree
than every other term of $r^3+xr+x^4+1$. If $\deg r=d\ge2$, the term
$r^3$ has degree $3d>\max\{d+1,4\}$ and cannot cancel. The zero polynomial
is not a root either. This is a contradiction in every case.

Consequently $f$ is irreducible over $\C(x)$ and over $\C[x,y]$: a factor
depending only on $x$ would have to divide the leading coefficient $1$ in
$y$, and any other factorization would persist over $\C(x)$.

\newpage
\subsection*{3. The projective curve is smooth}
The projective completion is
\[
  C:\quad F(X,Y,Z)=X^4+XYZ^2+Y^3Z+Z^4=0
  \quad\text{in }\mathbb P^2.
\]
Its defining polynomial is irreducible over $\C$: a projective
factorization would dehomogenize to one of $f$, unless a factor were a power
of $Z$, and $Z$ does not divide $F$.

\textit{Affine points.} A singular point with $Z\ne0$ would give complex
numbers $x,y$ satisfying
\[
  f(x,y)=0,\qquad 4x^3+y=0,\qquad x+3y^2=0.
\]
Multiply the last two equations by $x$ and $y$, respectively. We obtain
\[
  4x^4+xy=0,\qquad xy+3y^3=0.
\]
Together with $x^4+xy+y^3+1=0$, these force
\[
  \boxed{\quad x^4=\frac35,\qquad y^3=\frac45,
  \qquad xy=-\frac{12}{5}.\quad}
\]
The first two equalities imply $|x|<1$ and $|y|<1$, even when $x,y$ are
complex. Hence $|xy|<1$. The third equality instead gives
$|xy|=12/5>1$. Thus there are no affine singularities.

\textit{The point at infinity.} Setting $Z=0$ gives $X^4=0$, so the only
point at infinity is $P=[0:1:0]$. Since
\[
  F_Z=2XYZ+Y^3+4Z^3,\qquad F_Z(P)=1,
\]
this point is nonsingular as well. Therefore $C$ is smooth everywhere.

\subsection*{4. Genus and the finiteness conclusion}
The genus formula for a smooth geometrically irreducible projective plane
curve of degree $d$ over a field of characteristic zero is
\[
  g=\frac{(d-1)(d-2)}2.
\]
All its hypotheses were checked above, so our quartic has $g=3$.

\begin{faltings}
If a smooth, projective, geometrically irreducible curve defined over a
number field has genus at least two, then it has only finitely many points
rational over that number field.
\end{faltings}

Apply this theorem to $C$ over $\Q$. Its rational point set $C(\Q)$ is
finite. Each rational solution of the original equation gives the distinct
point $[x:y:1]\in C(\Q)$. Hence there are finitely many rational solutions,
and therefore finitely many integer solutions. Finally,
\[
  0^4+0(-1)+(-1)^3+1=0,
  \qquad 1^4+1(-1)+(-1)^3+1=0,
\]
which verifies the two displayed solutions and completes the proof.
\hfill$\square$

\textbf{What is established.} Infinitely many integer solutions are ruled
out unconditionally. This argument supplies neither a height bound nor a
complete list of integer solutions; in particular, it does not claim that
$(0,-1)$ and $(1,-1)$ are the only ones. Faltings' theorem and the plane-curve
genus formula are the explicitly identified external mathematical inputs.
\end{document}


